% ============================================================================
% CHAPTER 5 - RESULT ANALYSIS   (template file chapter_6.tex)
% Rubric mapping: CO6 (10 marks, 5 panel + 5 supervisor) is assessed from
% 5.1, 5.2, 5.3, 5.4, 5.5 and 5.6. CO7 (5 marks) is assessed from 5.3.
%
% EVERY number in this chapter comes from RESULTS.md in the project repository,
% where each one has the command that regenerates it.
% ============================================================================

\section{Performance Evaluation}

\subsection{What was measured, and the headline figures}

Each stage of the pipeline is evaluated on its own, because a single end-to-end
number would hide which part works. Table \ref{tab:headline} collects the main
results, and the rest of this section takes them one at a time.

\begin{table}[H]
\centering
\caption{The main results, one per stage. Each row has its own section below,
where the conditions and the caveats are stated.}
\label{tab:headline}
\small
\begin{tabular}{p{4.0cm}p{4.4cm}p{5.0cm}}
\toprule
\textbf{Stage} & \textbf{Measure} & \textbf{Result} \\
\midrule
Speech, six unseen lectures &
Median per-clip character error rate &
\textbf{67.7\% to 15.8\%} and 16.0\% in two seeds \\
\addlinespace
Speech, unseen lecturer &
Median per-clip character error rate &
72.8\% to 50.3\%, all six runs significant \\
\addlinespace
Board reconstruction &
Tiles finding a clear view, 35 boards &
median 97.7\% \\
\addlinespace
Board reconstruction &
Human check, 98 boards &
82 complete, 9 with writing genuinely lost \\
\addlinespace
Board reading &
Recall of 349 hand-verified items &
\textbf{31.2\% to 88.8\%} by changing the prompt \\
\addlinespace
Generated notes &
Recall of 349 hand-verified items &
\textbf{37.2\% to 89.1\%} \\
\bottomrule
\end{tabular}
\end{table}

\subsection{Speech recognition: the final result}
\label{sec:asr-result}

The corpus was closed at 28 transcripts and 5.15 hours. Whole lectures were held
out at random, 20 per cent of each lecturer's minutes with the seed fixed at zero,
giving 22 training lectures totalling 4.10 hours and six test lectures totalling
1.05 hours and 177 clips. The three validation lectures used for tuning were locked
out of the test set. Table \ref{tab:asrfinal} gives the result.

\begin{table}[H]
\centering
\caption{Median per-clip error on the 177 test clips, off-the-shelf against
fine-tuned. Greedy decoding throughout. The last row is the same configuration as
the row above it with a different random seed.}
\label{tab:asrfinal}
\small
\begin{tabular}{p{2.6cm}p{1.2cm}p{1.4cm}p{1.4cm}p{2.0cm}p{1.8cm}p{1.8cm}}
\toprule
\textbf{Learning rate} & \textbf{Seed} & \textbf{CER} & \textbf{WER} &
\textbf{Runaway clips} & \textbf{Better of 177} & \textbf{Wilcoxon} \\
\midrule
off-the-shelf & -- & 67.7\% & 93.9\% & 13 & -- & -- \\
\textbf{1e-3} & \textbf{42} & \textbf{15.8\%} & \textbf{42.5\%} & 1 & 164 &
p $<$ 1e-30 \\
\textbf{1e-3} & \textbf{1} & \textbf{16.0\%} & \textbf{41.7\%} & 1 & 167 &
p $<$ 1e-30 \\
2e-3 (the tuned choice) & 1 & 16.2\% & 41.7\% & 2 & 168 & p $<$ 1e-30 \\
2e-3 (the tuned choice) & 42 & \textbf{118.7\%} & 189.8\% & \textbf{82} & 35 &
worse \\
\bottomrule
\end{tabular}
\end{table}

The headline is that the character error rate falls from 67.7 per cent to 15.8 and
16.0 per cent, and the word error rate from 93.9 per cent to 42.5 and 41.7 per
cent. Two seeds agree to within 0.2 points, and the fine-tuned model is better on
164 and 167 of the 177 clips. About three quarters of the character error is
removed.

\begin{figure}[H]
\centering
\includegraphics[width=0.9\textwidth]{fig-asr-final}
\caption{The final speech result. Median per-clip error before and after
fine-tuning, for both stable seeds, with the diverged seed shown beside them
rather than left out.}
\label{fig:fig-asr-final}
\end{figure}

\begin{figure}[H]
\centering
\includegraphics[width=0.9\textwidth]{fig-asr-distribution}
\caption{The distribution of per-clip character error rate over the 177 test
clips, before and after fine-tuning. The shift is of the whole distribution, not
of a few easy clips.}
\label{fig:fig-asr-distribution}
\end{figure}

\begin{figure}[H]
\centering
\includegraphics[width=0.9\textwidth]{fig-asr-clip-scatter}
\caption{Each point is one test clip, with its error before fine-tuning on one
axis and after on the other. Points below the diagonal improved. The clips at the
extreme right are the ones where the off-the-shelf model fell into a repetition
loop.}
\label{fig:fig-asr-clip-scatter}
\end{figure}

\textbf{The instability, which has to be reported with the result.} At 2e-3, which
is the learning rate the pre-registered procedure chose on the validation
lectures, one of the two seeds collapsed. It produced runaway output on 82 of the
177 clips and a character error rate of 118.7 per cent, which is worse than doing
nothing at all. The data, the split and every other setting were identical to the
seed that reached 16.2 per cent. The 16.2 per cent figure is therefore never quoted
in this thesis without its failed twin.

This is not the first time this configuration has shown the same fragility. LoRA
rank 32 diverged and adapting all four attention projections diverged, in both the
rehearsal and the full tuning run, which is reported in the next section.

\textbf{A disclosure about the 1e-3 rows.} The pre-registered procedure selected
2e-3. The runs at 1e-3 were launched after the divergence at 2e-3 had been seen, so
that choice was not blind, and the thesis says so rather than presenting 1e-3 as
the pre-registered winner. What can be claimed honestly is narrower and is still
worth claiming. The runner-up rate from the same tuning table is stable across two
seeds on the test set and its two seeds agree to 0.2 points, while the chosen rate
is not stable. The selection rule itself was not changed and remains recorded in
the committed tuning plan.

\subsection{How the training settings were chosen}
\label{sec:asr-tuning}

The tuning procedure of Section 4.4.2 was run on the full corpus, scored only on
the three validation lectures, and never on the test lectures. Off-the-shelf
Whisper on those validation clips gives a character error rate of 73.8 per cent and
a word error rate of 97.6 per cent. Table \ref{tab:tuning} is the result.

\begin{table}[H]
\centering
\caption{Pre-registered hyperparameter tuning on the full corpus, scored on the
three validation lectures. Each stage starts from the winner of the previous one.
The rows marked as diverged produced output worse than the off-the-shelf model.}
\label{tab:tuning}
\small
\begin{tabular}{p{2.9cm}p{2.3cm}p{1.1cm}p{1.0cm}p{1.4cm}p{1.1cm}p{1.2cm}p{1.2cm}}
\toprule
\textbf{Stage} & \textbf{Run} & \textbf{lr} & \textbf{rank} & \textbf{layers} &
\textbf{epochs} & \textbf{val CER} & \textbf{val WER} \\
\midrule
1 learning rate & lr 5e-4 & 5e-4 & 16 & q,v & 8 & 37.9\% & 60.0\% \\
1 learning rate & lr 1e-3 (default) & 1e-3 & 16 & q,v & 8 & 42.2\% & 61.9\% \\
1 learning rate & \textbf{lr 2e-3} & 2e-3 & 16 & q,v & 8 & \textbf{33.6\%} &
\textbf{54.5\%} \\
2 LoRA rank & rank 8 & 2e-3 & 8 & q,v & 8 & 40.1\% & 64.1\% \\
2 LoRA rank & rank 32 & 2e-3 & 32 & q,v & 8 & \textbf{126.6\% diverged} &
200.0\% \\
3 adapted layers & all four projections & 2e-3 & 16 & q,k,v,o & 8 &
\textbf{134.7\% diverged} & 174.4\% \\
4 epochs & 5 epochs & 2e-3 & 16 & q,v & 5 & 39.4\% & 59.6\% \\
5 stability & seed 1 & 2e-3 & 16 & q,v & 8 & 39.1\% & 59.6\% \\
\bottomrule
\end{tabular}
\end{table}

The procedure chose a learning rate of 2e-3, rank 16, the query and value
projections and 8 epochs. That configuration is 8.6 character error points better
than the starting point, and the acceptance rule required it to beat the default by
more than the spread between two seeds, which here is 5.4 points.

Three findings come out of this table, and none of them is flattering.

\textbf{More adapter capacity hurts at this learning rate.} Rank 32 diverged
outright, and adapting all four attention projections diverged as well. More
trainable parameters are not better here. Five hours of speech does not support
them.

\textbf{Validation loss is a poor guide to the metric that is reported.} The
validation loss is lowest at 5 epochs, at 1.666 against 1.801 at 8 epochs, so stage
4 retrained there. The character error rate got worse, from 33.6 to 39.4 per cent.
The lesson, which is worth stating because it is a common mistake, is to select on
the metric that will be reported rather than on the loss.

\textbf{The margin is not large next to seed noise.} The winner is 8.6 points ahead
of the default with a 5.4 point spread from a single seed change, measured on 34
minutes of validation speech. It passes the rule that was written down in advance,
and it is still a modest effect. The outcome on the test set, where that same
configuration diverged on one seed, is consistent with a genuinely noisy region of
the setting space rather than with a clear winner.

A rehearsal of the same procedure had been run earlier on 2.69 hours of transcript,
while transcription was still going on, to check that the machinery worked. It
chose the same learning rate, and it produced the same two divergences at rank 32
and at four projections. Its full table is in the project record.

\subsection{Decoding, and the loop safeguard}
\label{sec:asr-safeguard}

An earlier result in this project, on a smaller model, only beat its baseline when
Whisper's compression-ratio loop safeguard was switched on. The obvious challenge to
the present result is therefore that the safeguard is doing the work. It is not.
Table \ref{tab:safeguard} rescored both final models with the safeguard on.

\begin{table}[H]
\centering
\caption{The final model scored with plain greedy decoding and with Whisper's loop
safeguard. The safeguard is applied identically to both models and reads only the
hypothesis.}
\label{tab:safeguard}
\small
\begin{tabular}{p{2.0cm}p{2.6cm}p{3.4cm}p{4.4cm}}
\toprule
\textbf{Seed} & \textbf{Greedy CER} & \textbf{With the safeguard} &
\textbf{Runaway clips, greedy} \\
\midrule
42 & 15.8\% & 15.8\% & 1 of 177 \\
1 & 16.0\% & 16.0\% & 1 of 177 \\
\bottomrule
\end{tabular}
\end{table}

The two are identical to the first decimal place, because the fine-tuned model loops
on one clip in 177 and the safeguard has nothing to rescue. Off-the-shelf Whisper
loops on thirteen. All figures reported in this thesis are the plain greedy ones.

\subsection{Where the gain is uneven}
\label{sec:asr-perlecture}

The improvement is not evenly spread, and the pattern follows the training data.
Table \ref{tab:perlecture} breaks the result down by test lecture.

\begin{table}[H]
\centering
\caption{Median character error rate per test lecture, seed 42, with the amount of
that lecturer's speech in the training set.}
\label{tab:perlecture}
\small
\begin{tabular}{p{2.8cm}p{2.0cm}p{1.4cm}p{2.4cm}p{2.0cm}p{2.6cm}}
\toprule
\textbf{Test lecture} & \textbf{Lecturer} & \textbf{Clips} &
\textbf{Off-the-shelf} & \textbf{Fine-tuned} &
\textbf{That lecturer in training} \\
\midrule
BanglaASR8 & A & 14 & 68\% & 15\% & 1.2 h \\
BanglaASR9 & A & 41 & 67\% & 12\% & 1.2 h \\
\textbf{BanglaASR11} & \textbf{B} & 46 & 70\% & \textbf{52\%} & \textbf{0.7 h} \\
BanglaASR15 & C & 8 & 88\% & 12\% & 2.2 h \\
BanglaASR19 & C & 12 & 65\% & 12\% & 2.2 h \\
BanglaASR27 & C & 56 & 73\% & 16\% & 2.2 h \\
\bottomrule
\end{tabular}
\end{table}

\begin{figure}[H]
\centering
\includegraphics[width=\textwidth]{fig-asr-per-lecture}
\caption{Median character error rate for each test lecture before and after
fine-tuning. Every lecture improves. Lecturer B's lecture improves least, and
lecturer B contributed the least training speech.}
\label{fig:fig-asr-per-lecture}
\end{figure}

Every lecture improves, but lecturer B's lecture improves by far the least, ending
at 52 per cent where the others reach 12 to 16 per cent, and lecturer B contributed
the least training speech at 0.7 hours. With one test lecture per lecturer at the
low end, this is an observation and not a controlled result. It cannot separate
``less data from this lecturer'' from ``this lecturer is harder to transcribe''.
The honest reading is that it is the most likely explanation and the clearest
argument for collecting more speech from under-represented lecturers.

\subsection{Is the remaining error only a spelling problem?}
\label{sec:asr-spelling}

Banglish has no standard spelling, so a reasonable challenge is that the word error
rate is inflated by \textit{aami} against \textit{ami}. Two spelling-tolerant
variants were computed on the same saved outputs, with the rules committed before
any rescoring, and are shown in Table \ref{tab:spelling}. These numbers come from
the six leave-one-speaker-out runs rather than from the final model, because that
was the headline at the time the rules were fixed.

\begin{table}[H]
\centering
\caption{Spelling-tolerant word error rates, mean of six per-run medians. nWER maps
both sides to the transcription guide's canonical spellings and collapses repeated
vowels. fWER additionally accepts a word at 80 per cent character similarity, so
words of three letters or fewer must still match exactly.}
\label{tab:spelling}
\small
\begin{tabular}{p{5.0cm}p{3.0cm}p{3.0cm}}
\toprule
\textbf{Measure} & \textbf{Base} & \textbf{Fine-tuned} \\
\midrule
Plain word error rate & 95.3\% & 75.7\% \\
Spelling-fair nWER & 95.0\% & 74.8\% \\
Fuzzy fWER & 92.7\% & 68.2\% \\
\bottomrule
\end{tabular}
\end{table}

The canonical spelling table moves the word error rate by about one point, and
accepting near-miss spellings moves it by about seven more. Most of the word error
is therefore genuine recognition error rather than spelling disagreement. The
fine-tuned model is better in all six runs on all three measures, with p between
2.2e-16 and 2.7e-06. Plain word and character error rates remain the headline, and
this table exists to answer the spelling question when it is asked.

\subsection{An unseen lecturer, and a correction to an earlier result}
\label{sec:asr-loso}

The result in Section \ref{sec:asr-result} holds out whole lectures from lecturers
who are also heard in training. A harder question is what happens on a lecturer the
model has never heard. Each lecturer was held out in turn and the model trained on
the other two, with two seeds per fold and plain greedy decoding.

\begin{table}[H]
\centering
\caption{Leave-one-speaker-out. Each row holds out one lecturer entirely.}
\label{tab:loso}
\small
\begin{tabular}{p{2.4cm}p{2.6cm}p{3.2cm}p{3.2cm}p{2.0cm}}
\toprule
\textbf{Held out} & \textbf{Trained on} &
\textbf{WER, base to tuned} & \textbf{CER, base to tuned} &
\textbf{Wilcoxon} \\
\midrule
A, 172 clips & B + C, 82.9 min & 97.5\% to 76.2 / 75.8\% &
74.2\% to 51.8 / 49.9\% & p $<$ 1e-10 \\
B, 184 clips & A + C, 80.0 min & 93.8\% to 73.5 / 74.8\% &
70.0\% to 48.8 / 52.1\% & p $<$ 1e-08 \\
C, 75 clips & A + B, 114.0 min & 94.6\% to 74.5 / 79.8\% &
74.2\% to 45.2 / 54.0\% & p $<$ 1e-05 \\
\bottomrule
\end{tabular}
\end{table}

Averaged over the six runs, the character error rate falls from 72.8 per cent to
50.3 per cent and the word error rate from 95.3 per cent to 75.7 per cent. All six
runs are significant under plain decoding, and runaway clips fall or stay level in
every fold.

What this supports is that fine-tuning on 80 to 114 minutes of transcribed Banglish
from two lecturers cuts the character error on a third, unheard lecturer by about a
third, for every one of the three lecturers and in both seeds. What it does not
support is that this is usable speech recognition, since a 50 per cent character
error rate is still very high, or that the finding holds across many speakers, since
three is a small number.

Both designs belong in the thesis and both are labelled every time they appear. The
result in Section \ref{sec:asr-result} is ``new lectures from known lecturers'' on
4.10 hours of training speech, and this one is ``an unseen lecturer'' on 80 to 114
minutes.

\textbf{A correction that has to be recorded.} An earlier version of this work
reported a word error rate falling from 96.1 to 81.8 per cent on a speaker never
seen in training. That claim was wrong, and the reason is instructive. The original
nine transcripts carried no lecturer label, so the lectures were labelled by hand,
and one lecture was attributed to the wrong lecturer. That lecture was in the
training set and its lecturer was the test speaker, so 47 clips of the supposedly
unseen speaker had been trained on. The error was found by computing speaker
embeddings, which separate the groups without ambiguity at 0.99 to 1.00 similarity
within a lecturer and 0.75 to 0.80 across. The result was recomputed on the
corrected split, the old number was retired, and the voice check is now a standing
part of the data preparation. This is also the reason the corpus statistics in
Chapter 4 report the lecturer of every recording from voice rather than from a file
name.

\textbf{A data quality check on the same result.} Nine of 156 timestamps in the
ground truth behind these runs carried stray spaces, so their text joined the
previous segment and ten test clips carried more words than their audio. Both models
were scored on the same clips, so the comparison is fair, but the effect was measured
rather than assumed. Without those ten clips the result is 72.2 per cent to 49.6 per
cent character error, against the published 72.8 to 50.3. The defect made the
published figure slightly pessimistic, so the published figure stands.

\subsection{Board reconstruction: coverage}
\label{sec:board-coverage}

The reconstruction is measured first by coverage, which is the fraction of tiles
that found a moment when the lecturer was not in front of them. Across the 35
boards of the first nine lectures, the median board has 97.7 per cent of its tiles
fully clear, the mean is 96.5 per cent, and the range is 87.5 to 100 per cent. Six
boards reach 100 per cent, 26 of the 35 reach 95 per cent or better, and only two
fall below 90 per cent. Over all 145 boards produced in the project the median is
100 per cent and the mean 98.4 per cent, with the worst board at 76.6 per cent.

Two things drive the weak boards. Short eras give a tile fewer chances to be seen
clear, and six eras are under eight frames. A lecturer who stands in one place is
the other cause, and the worst lecture in the set is the one where the lecturer holds
a single position far longer than the others.

That second cause was addressed rather than excused. For the lecturer who moves
least, extracting frames every 2 seconds instead of every 10 gives each tile five
times as many chances. The number of his boards reaching 95 per cent clear rose from
1 to 5 out of 10, and the worst board improved from 59 to 88 per cent.

Coverage is a measure of the method, not of the result, and it is important not to
read it as fidelity. A tile can be clear and the writing can still be missing,
because the writing may never have been visible in any frame. That question needs a
person, and it is the next section.

\subsection{Board reconstruction: what a human check found}
\label{sec:board-human}

Every one of the 98 boards of the newer lectures was checked by hand. A page showed
each reconstructed board beside two real frames from the same period of the video,
and the question was whether the board contains everything that was written. Table
\ref{tab:completeness} gives the answers.

\begin{table}[H]
\centering
\caption{Human check of 98 reconstructed boards.}
\label{tab:completeness}
\small
\begin{tabular}{p{6.0cm}p{2.0cm}p{5.4cm}}
\toprule
\textbf{Verdict} & \textbf{Boards} & \textbf{} \\
\midrule
\textbf{Complete} & \textbf{82} & 83.7 per cent of 98 \\
Something missing & 16 & \\
\quad writing genuinely lost & 9 & 9.2 per cent \\
\quad era cut while the board was being wiped & 4 &
the board is half erased in every frame of that era \\
\quad camera out of focus for the whole era & 3 &
unreadable in the source video, not a reconstruction fault \\
\bottomrule
\end{tabular}
\end{table}

\begin{figure}[H]
\centering
\includegraphics[width=0.9\textwidth]{fig-board-completeness}
\caption{The human check of 98 reconstructed boards, with the causes of the
sixteen incomplete ones separated.}
\label{fig:fig-board-completeness}
\end{figure}

The nine real losses are small. Seven of them are a single token or one diagram,
such as a label \texttt{t4}, the number \texttt{100} or the word Hybrid. One board
is the bad case, where a whole subnetting tree is absent. Counting only the boards
whose source frames were usable, which means dropping the four mid-erase eras and
the three out-of-focus ones, the reconstruction is complete on 82 of 91 boards, or
90.1 per cent.

The four mid-erase cases are worth separating because they have a different fix.
They are not a mosaic fault at all. The era boundary was placed while the lecturer
was still wiping, so every frame in that era shows a half-erased board. Holding the
era open until the wipe finishes would fix it, and it was not attempted before
submission.

These numbers count different things from the coverage figures above and from the
recall figures below, and they should not be added together. Coverage counts tiles,
this counts whether writing survived, and recall counts whether a model could read
the writing back.

\subsection{Reading the board with a vision-language model}
\label{sec:board-reading}

The same Qwen2.5-VL-7B model was run twice on the same board images, changing only
the prompt. The first asks for keywords, which is what the original pipeline did.
The second asks for a full transcription of everything on the board, with every
number exactly as written and \texttt{[illegible]} rather than a guess. Table
\ref{tab:vlmread} gives the result against the hand-verified answer keys.

\begin{table}[H]
\centering
\caption{Board-content recall, same model, same images, only the prompt changes.}
\label{tab:vlmread}
\small
\begin{tabular}{p{4.6cm}p{2.2cm}p{2.8cm}p{4.0cm}}
\toprule
\textbf{Boards} & \textbf{Keyword prompt} & \textbf{Full transcription} &
\textbf{Keyword to transcription} \\
\midrule
10 boards, 194 items & 41.2\% & 96.9\% & better on 9, worse on 0, sign p = 0.0039 \\
\textbf{35 boards, 349 items} & \textbf{31.2\%} & \textbf{88.8\%} &
better on 34, worse on 0, sign p = 1.2e-10 \\
A third lecturer, 10 boards, 87 items & not run & 89.7\% & -- \\
\bottomrule
\end{tabular}
\end{table}

\begin{figure}[H]
\centering
\includegraphics[width=0.9\textwidth]{fig-board-reading}
\caption{Board-content recall by item kind, keyword prompt against full
transcription, over 35 boards and 349 items. The numbers written on the board go
from none found to almost all found.}
\label{fig:fig-board-reading}
\end{figure}

Broken down by kind of item over the 35 boards, the change is from 0 of 67 numbers
to 65 of 67, from 43 of 47 names to 45 of 47, from 14 of 111 pieces of code to 90 of
111, from 47 of 94 terms to 89 of 94, and from 5 of 30 phrases to 21 of 30. The
keyword prompt found not one number on any board. This is the clearest single result
in the vision half of the thesis, and it is a result about prompting rather than
about model capability.

\textbf{The prompt was chosen on the same boards this is reported on, and that has
to be said.} There was no held-out set when the choice between the two prompts was
made. To give a figure that selection cannot explain, the boards were split by
lecture using a rule fixed before any per-half number was looked at, with the odd
numbered lectures as the development half and the even numbered ones as the test
half. The split is by lecture rather than by board because boards within a lecture
share a lecturer, a topic, a camera and a whiteboard.

\begin{table}[H]
\centering
\caption{The prompt comparison re-checked on a held-out half of the boards.}
\label{tab:promptsplit}
\small
\begin{tabular}{p{3.6cm}p{1.6cm}p{1.4cm}p{1.8cm}p{2.4cm}p{2.2cm}}
\toprule
\textbf{Half} & \textbf{Boards} & \textbf{Items} & \textbf{Keyword} &
\textbf{Full transcription} & \textbf{Better / worse} \\
\midrule
development & 17 & 173 & 34.1\% & 93.1\% & 17 / 0 \\
\textbf{test} & \textbf{18} & \textbf{176} & \textbf{28.4\%} & \textbf{84.7\%} &
\textbf{17 / 0} \\
all, the headline & 35 & 349 & 31.2\% & 88.8\% & 34 / 0 \\
\bottomrule
\end{tabular}
\end{table}

On the test half, which played no part in the choice, the full transcription prompt
still gives 28.4 per cent to 84.7 per cent, better on 17 of 18 boards and worse on
none. The honest limit of this check is that it is a split applied after the original
comparison, not a protocol registered before it, so it does not undo the original
selection. It answers the narrower question of whether the conclusion survives on
boards that played no part in choosing the prompt, and it does.

\subsection{The generated notes}
\label{sec:notes-result}

The notes were built for all 13 lectures that have hand-verified board answer keys,
in both languages, always from a transcript produced by an adapter that never heard
that lecture's lecturer. Table \ref{tab:notes} gives the recall of board content.

\begin{table}[H]
\centering
\caption{Board-content recall of the generated notes, 35 boards of the first nine
lectures, 349 items.}
\label{tab:notes}
\small
\begin{tabular}{p{6.0cm}p{2.0cm}p{5.4cm}}
\toprule
\textbf{Notes} & \textbf{Recall} & \textbf{Against the original pipeline} \\
\midrule
Original pipeline & 37.2\% & -- \\
\textbf{Annotated notes, Banglish} & \textbf{89.1\%} &
better on 32 boards, worse on 0, sign p = 4.7e-10 \\
Annotated notes, English written directly & 55.9\% &
better on 29, worse on 4, sign p = 1.1e-05 \\
Annotated notes, English translated from the Banglish & 89.4\% & -- \\
Pasting the raw board transcription into the notes & 88.0\% &
better on 30, worse on 1 \\
\bottomrule
\end{tabular}
\end{table}

\begin{figure}[H]
\centering
\includegraphics[width=0.9\textwidth]{fig-notes-recall}
\caption{How much of the board reaches the notes, for the original pipeline and for
the annotated notes in each language and route.}
\label{fig:fig-notes-recall}
\end{figure}

Over all 13 lectures and 436 items the Banglish notes reach 87.2 per cent and the
translated English notes 87.4 per cent. On the third lecturer's ten boards, for which
no earlier notes exist, the Banglish notes reach 79.3 per cent.

The comparison that matters most is with the row at the bottom of the table. Pasting
the vision-language model's raw board transcription into the notes reaches 88.0 per
cent, and the annotated notes reach 89.1 per cent, which a paired test cannot
separate, at 14 boards better, 8 worse and p = 0.29. That is the point rather than a
disappointment. The paste reaches its number by being a paste. The annotated notes
reach the same number while being a readable document with numbered boxes,
quotations and a step by step explanation.

\textbf{The quotation checks.} Across the 13 Banglish files the checker kept 31
quotations that appear word for word in the transcript and deleted 8 that do not, a
rejection rate of 21 per cent. The English files contain 27 quotations given as
translations, and the word-for-word check does not run on them, because a
translation cannot be matched against a Banglish transcript. The supportable claim is
therefore specific. Quotations in the Banglish notes are verified word for word
against the transcript, and quotations in the English notes are model translations
and are not verified. The claim ``every quote in the notes is verified'' would be
false and is not made anywhere in this thesis. One reference to a box that does not
exist survived across the 13 Banglish files and is counted in the output rather than
hidden.

\section{Analysis of Design Solutions}

This section analyses the design choices against the alternatives, and reports the
measurements that decided them. Four of the measurements below are negative results.

\subsection{Three approaches that did not work}

\textbf{Dual recogniser fusion.} An English recogniser and a Bengali recogniser were
run on the same nine lectures and their outputs merged. The merged output scores 73.9
per cent on the term measure against the baseline's 73.2 per cent, a difference of
0.7 percentage points, with an exact paired permutation test at p = 0.32 and a
Wilcoxon signed-rank test at p = 0.36. Four of the nine lectures get worse. The word
error rate of the merged output is 146.8 per cent against the baseline's 82.6 per
cent. The conclusion is that fusion does not fix Banglish recognition, which is what
motivated fine-tuning instead.

\textbf{Biasing the decoder towards board words.} A sweep of the bias strength from
0 to 2.0 on one lecture found term recall of 8.8 per cent with the bias off and 4.1
per cent at every non-zero setting, with the transcript truncated from 279 to 149
characters. Every non-zero strength produced byte-identical output. The shape of that
result is the finding. This is not a tuning problem with a better setting waiting to
be found, because the mechanism fails as soon as it engages. The recorded output of
the sweep gives the optimal bias as zero, so the experiment concluded against itself.

\textbf{Occlusion as a pointer.} Where the lecturer stands might indicate what the
lecture is currently about, which would give a free alignment between speech and
board regions. The test needs no manual labels. If standing in front of a region
means working on it, that region should gain more ink after the lecturer leaves than
the regions he was not standing in front of. Over 191 occlusion episodes in nine
lectures, the occluded region gained more ink in 101 cases and an unoccluded region
gained more in 81, with a Wilcoxon test at p = 0.054 and an exact sign test at p =
0.159. This does not support the claim and is reported as unsupported. The most
likely reason is that the test measures writing while a lecturer standing in front of
a region is often explaining it and adding no ink at all, which is exactly the case
where the signal would have been most useful.

\subsection{The display clean-up: a measured negative}
\label{sec:board-cleanup}

The clean-up described in Section 4.4.3 makes every board a clean white image and
was judged by eye to be better. Read by the same vision-language model with the same
prompt against the same answer keys, it is slightly worse. Table \ref{tab:cleanup}
gives the comparison.

\begin{table}[H]
\centering
\caption{The vision-language model reading the mosaic board against the cleaned
board.}
\label{tab:cleanup}
\small
\begin{tabular}{p{3.6cm}p{1.6cm}p{2.6cm}p{2.8cm}p{2.0cm}p{1.6cm}}
\toprule
\textbf{Set} & \textbf{Boards} & \textbf{Mosaic} & \textbf{Cleaned board} &
\textbf{Better / worse} & \textbf{Sign} \\
\midrule
Lectures 1 to 9 & 35 & 334/349 (95.7\%) & 329/349 (94.3\%) & 0 / 5 & p = 0.0625 \\
Third lecturer & 10 & 74/87 (85.1\%) & 71/87 (81.6\%) & 1 / 2 & p = 1.0 \\
\textbf{Combined} & \textbf{45} & \textbf{408/436 (93.6\%)} &
\textbf{400/436 (91.7\%)} & \textbf{1 / 7} & \textbf{p = 0.0703} \\
\bottomrule
\end{tabular}
\end{table}

The difference is not significant either way, so the honest statement is that there
is no measured difference, with the point estimate against the clean-up. It is not
that the clean-up is worse.

Reading the changes item by item explains the shape. The one gain is on the board the
clean-up was built for, where an end-of-era diagram that the old mask lost becomes
readable, and that board rises from 57.1 to 71.4 per cent. The losses are all single
fine details of the kind a clean-up erases, which are long code strings, a written out
sentence and two decimal numbers. There is also very little room to win, since the
mosaic is already at 95.7 per cent on the first nine lectures and 37 of the 45 boards
score identically on both.

The conclusion drawn is therefore narrow and is followed in the rest of the thesis.
The cleaned boards are kept because of how they look in a note, where every board is a
clean white image, and not because they help the model read. No claim is made that the
rebuild improved board reading, and the headline board numbers are reported on the
mosaic boards.

\subsection{Does the result need a large vision-language model?}
\label{sec:vlm-size}

The board reading was run again with Qwen2.5-VL-3B, changing nothing but the model.

\begin{table}[H]
\centering
\caption{Board reading with a 3B model against a 7B model, same prompt, same
images, same keys.}
\label{tab:vlmsize}
\small
\begin{tabular}{p{4.6cm}p{2.4cm}p{2.4cm}p{2.4cm}p{2.0cm}}
\toprule
\textbf{Boards} & \textbf{Qwen2.5-VL-3B} & \textbf{Qwen2.5-VL-7B} &
\textbf{7B better / worse} & \textbf{Sign} \\
\midrule
35 boards, 349 items & 85.4\% & 88.8\% & 11 / 5 & p = 0.21 \\
Third lecturer, 10 boards, 87 items & 77.0\% & 89.7\% & 5 / 1 & p = 0.22 \\
\textbf{Combined, 45 boards, 436 items} & \textbf{83.7\%} & \textbf{89.0\%} &
\textbf{16 / 6} & \textbf{p = 0.053} \\
\bottomrule
\end{tabular}
\end{table}

Changing the prompt is worth 57.6 percentage points, from 31.2 to 88.8. Changing the
model from 3B to 7B is worth 5.3 points, and that difference is not significant by the
sign test on either set on its own. A 3B model at 85.4 per cent already carries most of
the board. The 7B is the better model and remains the one reported, and the thesis does
not present model scale as the thing that made board reading work. This agrees with the
published finding that large multimodal models are strongly sensitive to how the
instruction is written \cite{ismithdeen2025promptception}.

The 32 billion and 72 billion parameter versions were not run, and the reason is
recorded rather than left as an omission. Their weights are 68.3 GB and 146.8 GB, the
machine had 62 GB of free disk at the time, and both would have had to be quantised to
fit in 32 GB of video memory, which would confound model size with quantisation loss.
This is therefore a comparison of a 3B model against a 7B model, and the thesis calls
it that rather than a scale study.

\subsection{Does the reconstruction help the model read the board?}

The reconstruction is a visual deliverable, and it is tempting to claim that it also
improves reading. The measurement does not support that. On the 35 boards of the first
nine lectures, recall rises from 88.8 per cent on a raw frame to 95.7 per cent on the
reconstructed board, better on 8 boards and worse on 3, with a sign test at p = 0.23
and a Wilcoxon test at p = 0.056. On the third lecturer's boards it goes the other way,
from 89.7 per cent to 85.1 per cent. This is a trend on one set and absent on the
other. The sentence this thesis uses is that the prompt is what makes the model read
the board, and the reconstruction is what makes the board usable in a note.

Two limits of the reconstruction were exposed by the hand check of the answer keys and
each deserves a sentence. Glare from the ceiling lights sits in every frame, so no
choice of frame removes it. Content that is visible in only one frame can be lost,
which happened to a binary number that was fully written at one moment, covered ten
seconds later and wiped ten seconds after that.

\subsection{Does fine-tuning the speech model improve the notes?}
\label{sec:2x2}

This is the question that joins the two halves of the thesis, and the answer is no, in
a way that needs explaining rather than hiding.

Four sets of notes were built for the 13 scored lectures, Banglish only, with
everything held constant except two things, which are whether the transcript came from
the fine-tuned model or the off-the-shelf model, and whether the board text was
supplied at all.

\begin{table}[H]
\centering
\caption{Board-content recall for the two-by-two, 35 boards, 349 items.}
\label{tab:2x2items}
\small
\begin{tabular}{p{4.6cm}p{4.0cm}p{4.0cm}}
\toprule
& \textbf{With board text} & \textbf{Without board text} \\
\midrule
\textbf{Fine-tuned transcript} & \textbf{89.1\% (311 items)} & 70.8\% (247) \\
Off-the-shelf transcript & 78.5\% (274) & 73.6\% (257) \\
\bottomrule
\end{tabular}
\end{table}

\begin{table}[H]
\centering
\caption{The same four cells counted by board, which is the test that matters.}
\label{tab:2x2boards}
\small
\begin{tabular}{p{5.6cm}p{2.2cm}p{2.6cm}p{1.6cm}p{1.6cm}}
\toprule
\textbf{Comparison} & \textbf{Item change} & \textbf{Better / worse} &
\textbf{Sign} & \textbf{Wilcoxon} \\
\midrule
With board text, base to fine-tuned & +10.6 pp & 7 / 6 & p = 1.0 & p = 0.81 \\
Without board text, base to fine-tuned & -2.8 pp & 10 / 9 & p = 1.0 & p = 0.72 \\
Fine-tuned, no board text to board text & +18.3 pp & 13 / 7 & p = 0.26 & p = 0.16 \\
Off-the-shelf, no board text to board text & +4.9 pp & 10 / 4 & p = 0.18 & p = 0.052 \\
\bottomrule
\end{tabular}
\end{table}

None of the four comparisons is significant. Read by board, the two speech conditions
split 7 to 6 with board text and 10 to 9 without it, which is as close to a coin toss
as this data can produce.

The eye-catching plus 10.6 points in the item totals is worth taking apart, because
quoting it as the benefit of fine-tuning would be wrong. Of the 37 items gained, two
boards of one lecture give 50 between them, and across the other 31 boards the
fine-tuned transcript is 13 items worse. One apparently large effect is two boards out
of thirty-five with the rest pointing gently the other way.

\textbf{Why the better transcript loses, and why that is a finding.} The two
transcripts were scored directly against the same answer keys, before any notes were
written.

\begin{table}[H]
\centering
\caption{How much of the board content each transcript already contains, 35 boards,
349 items.}
\label{tab:metricbias}
\small
\begin{tabular}{p{7.0cm}p{6.4cm}}
\toprule
\textbf{Transcript} & \textbf{Board items it contains} \\
\midrule
Off-the-shelf Whisper, character error rate 67.7\% & \textbf{106 (30.4\%)} \\
Fine-tuned Whisper, character error rate 15.8\% & 91 (26.1\%), better on 1 board,
worse on 9, sign p = 0.021 \\
\bottomrule
\end{tabular}
\end{table}

The far worse transcript contains significantly more of the board's content, and the
mechanism is mechanical rather than mysterious. Off-the-shelf Whisper translates into
English, and the answer keys are largely English technical strings such as NAND gate,
MySQL and lines of SQL. Its output is unusable as a transcript while still emitting
those exact strings. The fine-tuned model writes what the lecturer actually said, in
Banglish, and \textit{amra NAND gate dekhbo} does not string-match an English key
item. This is the same trap as the term measure discussed in Section 5.4.3. A measure
built out of English strings rewards a model that translates and penalises one that
transcribes.

\textbf{What the null result does and does not say.} It says that the fine-tuned
transcript changes almost nothing that board-content recall can see, and that measure
asks only whether facts written on the board reach the notes. It does not say the
transcript is irrelevant to note quality. The measure is blind by construction to
whether the explanation around those facts is correct, readable or even in the right
language, and the two transcripts differ enormously as transcripts. An English page
built from a transcript with 67 per cent character error will contain sentences that
are simply wrong, and nothing here measures that. Testing it properly needs readers,
which is why the reader study is the first item of future work. Until that exists, the
defensible claim is the narrow one, which is that the factual content of the notes
comes from the vision side and the speech improvement shows up in the quotations and
in the readability rather than in board recall.

\subsection{Two routes to the English notes}

Writing the English note directly from the board and the transcript gives 55.9 per
cent recall. Writing the section in Banglish first and translating it gives 89.4 per
cent, almost exactly the Banglish original's 89.1 per cent. That supports the reading
above, which is that the deficit of the direct route is the model paraphrasing the
board's strings into English prose rather than the note carrying less of the lecture.

The caveat is large and is stated with the number every time. Of the 130 items gained,
113, or 87 per cent, come from four dense boards carrying 39 to 42 items each, and
across the other 31 boards the two routes are level. Counted by board the comparison is
10 better, 7 worse and 18 tied, with a sign test at p = 0.63. There are real
regressions, with two boards falling from 6 of 7 items to 2 of 7. The statement used in
this thesis is therefore that the translated route is much better on dense,
number-heavy boards and indistinguishable from the direct route on ordinary ones.

\section{Final Design Adjustments}
\label{sec:adjustments}

Evaluation changed the design in nine places. This section lists what was changed and
what evidence forced it, and then states what is still wrong.

\subsection{Adjustments made because of a measurement}

\textbf{The loop safeguard was adopted, and then shown to be unnecessary.} Under plain
greedy decoding the small model fell into repetition loops on 58 of 184 clips, and the
loops decided the result. Whisper's compression-ratio safeguard was adopted, applied
identically to both models, and reported beside the plain numbers. With the large
model the fine-tuned system loops on one clip in 177, so the safeguard changes nothing
and the reported figures are the plain ones. Both facts are reported, including that
the safeguard was adopted after the greedy result had been seen.

\textbf{The person mask was replaced.} The original mask took whatever differed from
the usual appearance of the board, which failed when the lecturer stood still and
failed again when the camera's automatic exposure darkened a frame. It was replaced by
a pretrained segmentation network plus a shadow mask, which recovers writing the old
mask lost.

\textbf{Frames are extracted more densely for a lecturer who does not move.} Measured
in Section \ref{sec:board-coverage}.

\textbf{Each language version is written wholly in its own language.} The first design
followed a mock-up in which an English note quoted the lecturer in Banglish with a
translation beside it. Reading the first real notes showed that an English-medium
reader cannot use a Banglish quotation, so the rule was changed. This has a
consequence for what can be claimed, and the consequence is stated rather than glossed
over. The word-for-word check runs on the Banglish notes and cannot run on the English
ones.

\textbf{The English note is produced by translating the Banglish note.} Selected on the
measurement in Section 5.2.6, with the caveat about the four dense boards.

\textbf{A loop guard in the board reader.} On one board the model emitted the same
backslash line 512 times, giving 42,037 characters. Repeated identical lines are now
collapsed and the number dropped is recorded. This changes no score, because repeated
backslashes match no answer key item. It matters because that text is pasted into the
note prompt.

\textbf{A salvage parser in the board reader.} Boards carrying code make the model emit
unescaped quotation marks that break the JSON it returns, and without a fallback such a
board loses every one of its boxes. The parser now reads the fields one at a time when
the JSON cannot be parsed.

\textbf{A placeholder strip in the note builder.} Three times across 26 files the model
copied the prompt's own example text into the notes, producing a sentence of the form
\textit{The lecturer mentioned, ``exact words from the transcript''}. The quote checker
could not catch it, because it looks for quotations absent from the transcript and this
is not a quotation at all. The strip removes the prompt's placeholder text and records
how many it removed. The affected lectures were rebuilt and recall was unchanged, so
this is a readability fix rather than a scoring one.

\textbf{Label repair in the translated notes.} The translation left 97 section labels in
Banglish across the thirteen English files and returned three summary lines in capital
letters. Both are repaired by a mapping that does not call the model. The capitalisation
repair takes each word's casing from how the same file writes that word elsewhere, so
that Python, TTL and DBMS are not lower-cased.

\textbf{Normalising the Background label.} The model wrote the heading of the
non-lecture box in four different ways, and 11 of 21 boxes in one rebuild used a bare
heading with no disclaimer. Since the heading is what the reader sees, those boxes were
fenced but not labelled, which defeats the purpose of the box. The page builder now
normalises the label at render time, and all 62 boxes across the 26 pages now carry the
disclaimer.

\subsection{What is still wrong}
\label{sec:known-defects}

Four defects are known, counted and not fixed, and they are listed here rather than
left for a reader to find.

\textbf{The box finder treats a sticker as content.} The whiteboard in one room carries
a manufacturer's sticker, and because it is dark and separate the box finder puts a box
around it on every board of that lecture. It is visible in Figure
\ref{fig:board-demo-5}. It is cosmetic, it costs nothing in any measurement, and it
would be fixed by ignoring a region that appears unchanged across every era of a
lecture.

\textbf{Some quotations cannot be translated and are dropped.} On the thirteen scored
lectures the translation route kept 15 quotations and dropped 19 because it could not
produce a faithful version, and over all 42 lectures that have notes it kept 60 and
dropped 49. A dropped quotation costs the page a lecturer's voice. It is counted in the
output file of every page rather than being silent, and it is the largest remaining
defect in the English version.

\textbf{A few references to boxes that do not exist survive.} Across the 13 scored
Banglish files the box reference check removed all but one. Over all 42 lectures ten
survive in each language. They are counted, not silent.

\textbf{Four era boundaries fall in the middle of an erase.} Described in Section
\ref{sec:board-human}. The fix is to hold the era open until the wipe finishes.

\subsection{A factual error in the generated prose, found by reading}

One further defect is not a software fault and cannot be fixed by a checker. In the
notes for one digital logic lecture the model wrote that there are two types of
universal gates, which are XOR gates. The lecturer and the board both say NAND and
NOR. Board-content recall cannot catch this, because it measures whether board items
appear in the notes and not whether the prose around them is true.

This is stated here deliberately. Nothing in this thesis measures the factual
correctness of the generated explanation, the planned reader study is what would, and
an example of the error is more useful to a reader than a claim that it does not
happen.

\section{Statistical Analysis}

\subsection{Which tests were used and why}

Every comparison in this thesis is paired, because both systems are run on exactly the
same clips or the same boards. Two tests are used on every comparison. The Wilcoxon
signed-rank test ranks the differences by size and asks whether positive and negative
differences are balanced \cite{wilcoxon1945}. The exact sign test only counts how many
items improved and how many got worse, which makes fewer assumptions and is the more
conservative of the two.

Both are reported wherever they disagree, and they do disagree in places. The
two-by-two of Section \ref{sec:2x2} has one comparison at Wilcoxon p = 0.052 and sign
p = 0.18, and neither is treated as significant.

\subsection{Medians, not means, and the direction of the difference}

Two rules were adopted after being caught out by their absence.

\textbf{Report medians.} Both the base and the fine-tuned model sometimes fall into a
repetition loop and emit far more words than were spoken, which pushes a mean error
rate above 100 per cent. On one early comparison the same data gave p = 0.76 on means
and p = 2.19e-05 on medians with a paired test. The mean result was an artefact of a
handful of runaway clips, not evidence of no effect.

\textbf{Read the direction, not only the p-value.} A two-sided p-value does not say
which system won. Under plain greedy decoding one early adapter was significantly worse
than its baseline, at z = -2.62 and p = 8.7e-03, and reporting only the p-value would
have made that look like a success. The sign of the test statistic is stated beside
every p-value in this thesis.

\subsection{A metric that moves the wrong way}

The measure inherited from the earlier phase of this project is a term-based F1 over a
fixed list of 82 English words. Two independent training runs on identical data, with
the same recipe and the same seed, differing only in the order the manifest listed the
clips, gave the results in Table \ref{tab:termf1}.

\begin{table}[H]
\centering
\caption{Two training runs on identical data, showing which measures replicate.}
\label{tab:termf1}
\small
\begin{tabular}{p{4.6cm}p{3.0cm}p{3.0cm}p{2.4cm}}
\toprule
\textbf{Measure} & \textbf{Run A} & \textbf{Run B} & \textbf{Spread} \\
\midrule
Word error rate, median & 81.8\% & 77.1\% & 4.7 pp \\
Character error rate, median & 60.7\% & 57.4\% & 3.3 pp \\
Term recall & 88.8\% & 88.3\% & 0.5 pp \\
Term precision & 66.7\% & 55.3\% & \textbf{11.4 pp} \\
Term F1 & 76.1\% & 68.0\% & \textbf{8.1 pp} \\
\bottomrule
\end{tabular}
\end{table}

The speech result replicates, with both runs beating the base model at p below 0.01 on
word and character error rate. Term F1 does not, swinging 8.1 points on identical data,
driven by an 11.4 point swing in term precision.

The reason is the interesting part. Term precision counts how many of the 82 fixed
English words the model emits. Run B transcribes better Banglish, so it emits fewer
English words, so its term precision falls. The measure penalises the model for doing
the task correctly. Three consequences follow and are applied throughout this thesis. A
term F1 difference smaller than about 8 points is never claimed, which independently
rules out the fusion result of 0.7 points. Word and character error rate with their rank
tests are the headline. Both runs are reported rather than the better one being quietly
chosen.

The same effect reappears in a different place, in the note evaluation of Section
\ref{sec:2x2}, where the worse transcript contains more of the English answer key
items. Two independent instances of the same failure in one project is the reason this
thesis treats it as a finding rather than as a nuisance, and it agrees with the
published work on evaluation metrics for code-switched recognition discussed in Section
2.2.4.

\subsection{Counting by board and counting by item}

Recall over answer key items is dominated by dense boards, because a board with 42
items contributes fourteen times as much as a board with 3. Two results in this thesis
look very different depending on which is counted. The translated English route gains
130 items over the direct route, of which 113 come from four boards, and counted by
board it is 10 better and 7 worse with p = 0.63. The two-by-two gains 37 items from
fine-tuning, of which 50 come from two boards while the other 31 boards lose 13, and
counted by board it is 7 better and 6 worse with p = 1.0.

Both countings are reported for every comparison where they differ, and the by-board
count is the one the conclusion rests on.

\subsection{Multiple comparisons}

This chapter reports a large number of tests. No correction for multiple comparisons
has been applied, and that is a limitation rather than an oversight. It matters least
for the speech result, where the p-values are below 1e-30 and would survive any
correction, and it matters most for the marginal results, which is why every marginal
result in this chapter is described as a trend or as not significant rather than being
counted as a finding.

\section{Comparisons and Relationships}

\subsection{Comparison with existing systems}

Table \ref{tab:neighbours} in Chapter 2 places this work beside the four nearest
Bengali-English code-switched speech systems. The comparison cannot be made on the
numbers, and it is worth saying why clearly. Two of the four are trained and evaluated
on synthetic speech, one targets the Bengali script and one is an undocumented public
model evaluated on a random split of its own clips. Error rates on different datasets
with different transcription targets are not comparable, and this thesis does not
present them as if they were.

What can be compared is the design of the evaluation, and there this work is stricter
than its neighbours in three ways. Whole lectures are held out, so no part of a test
recording appears in training. The references are checked word by word by a person
rather than being generated. The same test is reported under two designs, one with the
test lecturers heard in training and one with an entirely unseen lecturer, so that the
easier and harder questions are separated.

For the board side, the comparison with earlier work is different in kind. Systems such
as those of Davila and Zanibbi \cite{davila2017whiteboard} and Kota et al.
\cite{kota2019generalized} extract and binarise handwriting from fixed-camera lecture
video and report component-level recall above 94 per cent on public datasets. This work
does not beat them at that task and does not attempt to. Its reconstruction is
engineering assembled from known parts. What is added is the step after it, which is
reading the reconstructed board into text with a vision-language model, scoring that
against hand-verified answer keys, and carrying numbered regions of the board into a
generated note.

\subsection{The relationship between training data and gain}

Two observations connect the amount of data to the result, and neither is a controlled
experiment.

Within the final result, the lecturer with the least training speech has by far the
worst test lecture, at 52 per cent character error against 12 to 16 per cent for the
others. Across designs, holding out an entire lecturer and training on 80 to 114
minutes gives 50.3 per cent character error, while holding out whole lectures from
known lecturers and training on 4.10 hours gives 15.8 per cent. The two differ in both
the amount of data and the difficulty of the question, so neither comparison isolates
the effect of data volume.

An earlier scaling curve on the smaller model, at 0.30, 0.60 and 0.93 hours, gave 84.6,
80.6 and 83.3 per cent word error, which is flat within the 4.5 point run-to-run noise
after the first 18 minutes and significant at every budget. A curve on the current
corpus was planned and was cancelled when the corpus closed, so the relationship
between hours and accuracy at this scale is not measured here and is listed as future
work.

\subsection{The relationship between the stages}

The clearest relationship in the results is that the factual content of a note comes
from the vision side and not from the speech side. Adding the board text moves note
recall by 18.3 points with the fine-tuned transcript and 4.9 points with the
off-the-shelf one, while changing the transcript with board text present moves nothing
that a by-board test can detect. That is not a defect of the speech work. It is a
statement about what board-content recall measures, and it is the reason the speech
result is reported on its own terms as a transcription result rather than being
justified through the notes.

\subsection{Comparison with the earlier phase of this project}

The earlier phase of the project reported a fusion improvement of 5.7 percentage
points with p = 0.003 and a large effect size, an ablation table, a bias sweep and a
set of per-video statistics. None of those numbers was produced by any code in the
project. They were re-derived where a derivation was possible and retired where it was
not. The measured fusion effect is 0.7 points at p = 0.32. The bias sweep is the one in
Section 5.2.1. Two of the ablation rows described modes the pipeline does not have. A
failure-mode breakdown of five neat percentages summing to 100 had no analysis behind
it and was removed rather than reconstructed.

This is recorded here for two reasons. It is the reason for the project rule that a
number needs a command, and a reader comparing this report with the earlier one is
entitled to know which numbers changed and why.

\section{Discussions}

\subsection{What the results mean}

The speech result is the strongest part of the thesis. Fine-tuning a large multilingual
model with a small adapter on 4.10 hours of hand-checked code-mixed classroom speech
cuts the character error rate from 67.7 to about 16 per cent on lectures the model has
never heard, with two seeds agreeing and almost every clip improving. The practical
meaning is that the output changes from something a transcriber cannot use to a draft
a transcriber can correct.

The board result is a result about prompting. The same model on the same images goes
from 31.2 to 88.8 per cent recall of what is written, purely on how it is asked, and a
model less than half the size reaches 85.4 per cent. Anyone building a system like this
should spend their effort on the prompt before spending it on a bigger model.

The note result is that board content can be made to reach the page. The annotated
notes carry 89.1 per cent of the board's items against 37.2 per cent for a
conventional summarisation prompt, and they do it while remaining a readable document
rather than a paste of extracted text.

The negative results matter as much. Fusion, decoder bias and occlusion-as-pointer were
all built and all measured, and none of them worked. The display clean-up looks better
and reads slightly worse. Fine-tuning the speech model does not move note recall.
Reporting these is what makes the positive results credible.

\subsection{Implications}

For code-mixed speech, the implication is that the data is the bottleneck and the
method is not exotic. A standard adapter on a standard model gives a large gain, and
what took the effort was five hours of carefully checked transcription and a split
design that does not leak.

For evaluation, the implication is stronger and is the part of this thesis most likely
to be useful to others. Measures assembled out of English strings systematically reward
a system that translates code-mixed speech and punish one that transcribes it. It
happened twice here in two unrelated measures, and it is consistent with the published
work on code-switched evaluation reviewed in Section 2.2.4. Anyone evaluating a
code-mixed system on an English-derived metric should check for it before believing the
result.

For classroom systems, the implication is that the whiteboard is worth the trouble.
Most of the facts a student needs from these lectures are on the board and not in the
speech, and a pipeline that reads the board recovers them.

\subsection{Limitations}

The limitations are listed together so that none has to be found by a reader.

\begin{enumerate}
  \item \textbf{No human evaluation of the notes.} The measure counts whether board
        content reached the page. It cannot see whether the explanation is correct,
        readable or useful, and a factual error in the prose has already been found by
        reading.
  \item \textbf{Three lecturers, one university, one language pair.} Nothing here shows
        that the result transfers to another institution, another subject area or
        another code-mixed pair.
  \item \textbf{5.15 hours of transcribed speech.} Small by the standards of speech
        recognition. The relationship between hours and accuracy at this scale is not
        measured.
  \item \textbf{Training is unstable at the selected learning rate.} One seed of the
        pre-registered configuration diverged completely. The stable configuration
        reported beside it was chosen after that was seen.
  \item \textbf{The board-reading prompt was chosen on the boards it is reported on.}
        A held-out half was checked afterwards and supports the conclusion, but it does
        not undo the selection.
  \item \textbf{Box finding is not evaluated.} No layout ground truth exists. It works
        well on sparse boards and coarsely on dense ones, and it draws a box around a
        sticker.
  \item \textbf{Quotations in the English notes are not verified.} They are
        translations, and the automatic check cannot apply to them.
  \item \textbf{No comparison against a neighbouring model on our own data.} Running a
        public Bengali-English model on our test clips would turn a discussion into a
        table, and it was not done before submission.
  \item \textbf{No correction for multiple comparisons.} Marginal results are described
        as trends for that reason.
  \item \textbf{The system is for recorded lectures.} Nothing here runs in real time.
\end{enumerate}

\subsection{Threats to validity}

The main threat to the speech result would be leakage between training and test. Three
things guard against it, which are whole-lecture splits, the voice check on every
recording, and the frozen record of which lectures were on each side of every split.
One leak did occur earlier in the project, it was found by the voice check, and the
affected claim was retired.

The main threat to the vision results is that the answer keys were drafted by reading
the reconstructed boards, which could favour the reconstruction over a raw frame. The
keys were then checked by a person against the board images without seeing any model
output, and the comparison between the reconstruction and the raw frame is reported as
a trend rather than a result, so the conclusion does not depend on it.

The main threat to the note results is that board-content recall is a narrow measure
that a system could game by pasting the board text into the page. That is exactly what
the variant at 88.0 per cent does, which is why the annotated notes are compared
against it rather than only against the original pipeline, and why the thesis says that
matching it while remaining readable is the claim.
